
\subsubsection{Interpretation of Results}

The 90\% experimental success rate demonstrates that testability is a predictable property via formal constraint-satisfiability verification in proof-of-concept settings. By validating predictions against post-hoc p-values rather than circular expert agreement, this work demonstrates non-circular ground truth for meta-research tool evaluation. The system outperforms random baseline (50\%) with high statistical significance (binomial p = 0.0002).

Cross-domain confound pattern transfer (93.33\% precision) reveals that documented confounds reflect fundamental deep learning design choices that generalize beyond modality boundaries. NLP confound patterns successfully flag vision confounds and training confounds, suggesting confound detection can leverage a shared pattern library rather than requiring domain-specific rule engineering for each modality.

The two null results (testable hypotheses yielding p = 0.679 and p = 0.757) validate the system's distinction between testability and guaranteed significance. Both hypotheses had valid (D,B,M) triples and no flagged confounds, yet interventions did not produce significant effects. This outcome confirms the conservative design: the system predicts experimental feasibility but does not guarantee effect sizes or statistical significance.

\subsubsection{Limitations}

\textbf{KB Coverage Ceiling (84\%):} Eight well-known datasets are missing from the January 2026 HuggingFace snapshot, causing false negatives where testable hypotheses referencing Pascal VOC, STL-10, SST-2, Common Voice, tieredImageNet, CUB-200, or Stanford Cars are incorrectly classified as ''not testable.'' This limitation reflects catalog incompleteness rather than extraction failures.

\textbf{Confound Pattern Incompleteness (15 patterns, 2017-2020):} The pattern database misses post-2020 confounds such as pre-training dataset confounds (transfer learning), contrastive loss confounds (self-supervised learning), and prompt template confounds (large language models). One false negative stems from this limitation.

\textbf{Keyword Extraction Brittleness (10\% recall):} Simple substring matching fails on paraphrased hypotheses where dataset/benchmark/metric names are expressed non-standardly. This brittleness causes false negatives where testable hypotheses with informal phrasing are missed.

\textbf{Proof-of-Concept Validation (Simplified Experiments):} Validation used mock data with known effect sizes to test verification pipeline logic rather than real-world experiments with actual datasets and full training runs. The 90\% success rate may overestimate real-world performance---anticipated degradation to 70-80\% due to dataset loading failures, training instabilities, and hyperparameter sensitivity in production settings.

\subsubsection{Future Work}

Multi-source KB aggregation for 95\%+ coverage. Living confound database with automated literature mining and crowdsourced contributions. Semantic (D,B,M) extraction via sentence-BERT for 70-80\% recall at <10\% false positive rate. Real-world experiment execution for external validity validation. Partial confidence scoring to expose uncertainty on borderline cases.
